\subsubsection{Etats physiques du carbone}
\label{subsubsection:relation_potentiel_chimique_CP-transition_phase}


\setcounter{equation}{0}
\setcounter{Q}{0}

Le carbon existe à l'état solide sous deux variétés allotropiques principales : le graphite et le diamant.

\Q Calculer les volumes molaires $V^\star_{m,i}$ (supposé indépendant de la pression) du carbone graphite et du carbone diamant.
\solution{On peut calculer le volume molaire à partir de la masse molaire du carbone et de la masse volumique du solide : 
\begin{align*}
V_m^\star (\text{C (graphite)} ) &= \frac{M( \ce{C} )}{\rho ( \text{graphite} )} \stackrel{AN}{=} 5.31 \cdot 10^{-6} ~ \si{\meter\cubed\per\mole} \\
V_m^\star (\text{C (diamant)} ) &= \frac{M( \ce{C} )}{\rho ( \text{diamant} )} \stackrel{AN}{=} 3.41 \cdot 10^{-6} ~ \si{\meter\cubed\per\mole} \\
\end{align*}
}

\Q Exprimer le potentiel chimique d'un corps pur condensé $\mu^\star (T,P)$ en fonction de postentiel chimique standard $\mu^{\star, \circ}(T)$, de son volume molaire $V^\star_{m}$, de la pression $P$ ainsi que de la pression standard $P^\circ$.
\solution{
\begin{align*}
V = \left( \frac{\partial G}{\partial P} \right)_{T} \implies V^\star_m = \left( \frac{\partial \mu^\star}{\partial P} \right)_{T}
\end{align*}
On intègre entre un état standard de variables d'état $(T, P^\circ)$ et un état quelconque de variables d'état $(T,P)$
\begin{align*}
\int_{\mu^\star (T, P^\circ)}^{\mu^\star (T, P)} d\mu' &= \int_{P^\circ}^P V^\star_m dP' \\
\mu^\star (T, P) &= \mu^{\star \circ} (T) +  V^\star_m (P - P^\circ)
\end{align*}
}

\Q \'{E}crire l'équation de réaction associée à la transition de phase diamant-graphite. Exprimer l'enthapie libre de réaction et en déduire une condition d'équilibre chimique sur les potentiels chimiques.
\solution{
\textbf{\'{E}quation de réaction : } \ce{C(graphite) <=> C(diamant)}
\textbf{Enthalpie libre de réaction : } \`{A} une température $T$ et une pression $P$ donnée :
\begin{align*}
\Delta_r G = \mu^\star_{grp} (T,P) - \mu^\star_{dia}(T,P) \stackrel{\implies}{eq} \mu^\star_{grp} (T,P) = \mu^\star_{dia} (T,P)
\end{align*}
}

\Q Sous quelle pression faut-il opérer à \SI{298}{\kelvin} pour préparer du carbone diamant à part du carbone graphite. 
\solution{
On peut développer l'expression pour trouver la pression $P$ d'équilibre des phases :
\begin{align*}
\mu^\star_{grp} (T,P) &= \mu^\star_{dia} (T,P) \\
\mu^{\star \circ}_{grp} (T) + V^\star_{m, grp} (P-P^\circ) &= \mu^{\star \circ}_{dia} (T) + V^\star_{m, dia} (P-P^\circ) \\
\end{align*}
d'où au final : 
\begin{align*}
P (T) = P^\circ + \frac{\mu^{\star \circ}_{dia} (T) - \mu^{\star \circ}_{grp} (T)}{V^\star_{m, grp} - V^\star_{m, dia}}
\end{align*}
\textbf{Application numérique :} $P = 1.51 \cdot 10^9 ~ \si{\pascal}$ à \SI{298}{\kelvin}
}


\paragraph*{Données}
\begin{itemize}
\item Masse molaire : $M(\ce{C}) = 12.01 ~ \si{\gram\per\mole}$
\item Potentiels chimiques standard et masses volumiques à \SI{298}{\kelvin} 
\begin{center}
\begin{tabular}{ccc}
\toprule
constituant & $\mu^{\star, \circ}$ & $\rho$ \\
& \si{\kilo\joule\per\mole} & \si{\kilo\gram\per\meter\cubed} \\
\midrule
graphite & 0 & 2260 \\
diamant & 2.87 & 3513 \\
\bottomrule
\end{tabular}
\end{center}
\end{itemize}